% The full chapter title is too wide for the running head, so the header mark is
% shortened once (see \judulpendek in preamble-en.tex). The chapter page and the
% table of contents keep the complete title. The macro must come *before* \chapter:
% \chaptermark consumes it, so placing it after leaks the short title to the next chapter.
\judulpendek{The Limits of Attack}
\chapter{The Limits of Attack: Obfuscated Gradients and NIDS Evasion}

A group of students added one line of rounding to the model input, then reported that their own
attack failed completely: the success rate fell from seventeen percent to zero point zero two
percent. That number went into the executive summary, then into the slides, then into the abstract.
Two weeks later an examiner asked for one thing: prove that the gradient the attacker used still
points in the right direction. It did not. Rounding is flat almost everywhere, so its derivative is
zero, so the attacker takes a zero step, so the attack that failed is really an attack that never
happened. When the examiner replaced the derivative with a reasonable approximation, the success
rate returned to seventeen percent within minutes.

This chapter is about that mistake. The first mistake is simple: rounding erases the gradient, and
erasing the gradient means blinding the attacker, not making the model strong. The second mistake is
more interesting: there are times when obfuscating the gradient really does raise the cost of
attack, so part of that robustness is real. Telling the two apart demands experiment, not rhetoric.

The second half of the chapter carries the same problem to a much noisier place: network intrusion
detection systems, or NIDS. There packets must still arrive, protocols must remain valid, decisions
are made in milliseconds, and the model is retrained every day. All of that makes an evasion attack
that is beautiful on paper impractical in the field. But ``impractical'' is not ``impossible''.

\section{Gradients that vanish, and the sense of safety born from them}
\tingkat{basic}

Chapter 29 taught two attacks that rest on one object: the gradient of the loss with respect to the
input. The one-step attack takes the direction of the sign of the gradient and steps by
$\varepsilon$,

\begin{equation}
x' = x + \varepsilon\,\operatorname{sign}\bigl(\nabla_x L(x,y)\bigr),
\end{equation}

whereas the many-step attack repeats a small step and projects the result back onto the
$\ell_\infty$ ball of radius $\varepsilon$,

\begin{equation}
x_{t+1} = \Pi_{B(x,\varepsilon)}\Bigl(x_t - \alpha\,\operatorname{sign}
\bigl(\nabla_x L(x_t,y)\bigr)\Bigr).
\end{equation}

Both work for one reason that already appeared in the chapter on linearization: for a small
perturbation the loss can be approximated by a linear function, so the direction that raises the loss
fastest is the direction of the gradient. If that gradient is correct, the attack works at the cost
of one backward pass. If that gradient is wrong, the attacker goes blind. Those two sentences are not
equivalent to ``if the gradient is wrong, the model is safe''. A blind attacker can still guess, and
a guess without a gradient, for example random noise on packets, can still succeed in some cases.
What is lost is only the cheapest road, not every road.

\index{obfuscated gradient}
The term \emph{obfuscated gradient} names the situation in which the gradient an attacker sees does
not give enough information about the direction that actually raises the loss, not because the model
is strong, but because of the way the model is wrapped. Athalye, Carlini, Wagner and colleagues gave
this symptom a sharper name: a false sense of security. In their review of nine defenses from a
machine learning conference, seven depended on an obfuscated gradient; six were broken completely and
one partially, at an attack cost roughly the same as an ordinary attack.

\begin{untukanak}
Imagine a teacher who does not want her students to cheat, so she seals the notebook shut with duct
tape. The honest student does have trouble. But the creative student does not stop: he notices that
the answer to question one is always larger than the answer to question two, then guesses the whole
sheet from that pattern. The tape does not erase the answers; it only blocks one road to them.

A model that rounds its input works like that tape. The numbers that enter the model are cut to a
fixed multiple, for instance two decimal places, so that small changes are invisible altogether. A
program that tries to find the ``raise it a little at a time'' direction will find that nothing
changes, then give up and report that the model is immune. That report is wrong, and not because the
program is stupid. The program simply asked in a way that does not fit.

Now the part that matters most to remember. The same model can still be wrong if the numbers are
changed enough to cross a cut boundary. So the right question is not ``does a small attack succeed'',
but ``how large a change is needed''. This chapter teaches how to ask that question, and how to
recognize a report that never asks it.
\end{untukanak}

Why does rounding erase the gradient? Because rounding is a step function. Take
$q>0$ and define $Q_q(t)=q\,\mathrm{round}(t/q)$, the nearest multiple of $q$. On
the interval $kq-q/2 \le t < kq+q/2$ this function is constant at $kq$, so its derivative is zero
inside the interval and undefined at the boundary points. The set of boundary points has measure
zero, so for any smooth function $f$, the composition $f\circ Q_q$ has zero gradient almost
everywhere.

\begin{proposisi}
Let $Q_q\colon\mathbb{R}^d\to\mathbb{R}^d$ act coordinate by coordinate as above, and let
$f\colon\mathbb{R}^d\to\mathbb{R}$ be continuously differentiable. Then
$\nabla (f\circ Q_q)(x)=0$ for almost every $x\in\mathbb{R}^d$.
\end{proposisi}

\begin{proof}
Partition $\mathbb{R}^d$ into half-open boxes of side $q$, bounded by the
hyperplanes $x_j=k_jq\pm q/2$. In the interior of every box, $Q_q$ is
constant, so $f\circ Q_q$ is also constant and its gradient is zero. The complement of the union of
all those interiors is the union of the boundary hyperplanes, which has measure zero. So the
gradient is zero almost everywhere.
\end{proof}

\begin{akibat}
A one-step attack that uses the exact gradient of $f\circ Q_q$ produces
$x'=x$ for every $\varepsilon>0$. The measurable success rate is therefore
exactly zero, and the reported ``accuracy under attack'' is precisely the clean accuracy.
\end{akibat}

This corollary is the heart of the whole chapter. The statement is not about the strength of the
model; it is about the arithmetic of the attacker, and the zero in the results table does not measure
robustness, but the absence of a step. If the ``accuracy under attack'' column is identical to the
clean accuracy down to two decimal places, you are not looking at a defense that worked, but at an
attacker that stopped.

\begin{bacakeras}
Two voices. The first carries a map, the second carries a flashlight.

--- Look. I have made the map I drew blurry. Every uphill road I erased, there are no arrows, no
numbers.
--- Then the mountain is gone.
--- No. The mountain is still there.
--- So what is gone is only the picture?
--- What is gone is only the picture.
--- Someone who can smell the soil can still climb.
--- Still can.
--- Someone who throws a pebble and listens to the echo can also still climb.
--- Also still can.
--- So your map is not a wall.
--- Your map is not a wall. Your map is only a map that refuses to answer.

Repeat the last line three times, slower and slower: a map that refuses to answer is not a wall. Not
a wall. Not a wall.
\end{bacakeras}

Athalye and colleagues divided the mechanism into three families. The first family, shattered
gradients, appears when part of the model is non-differentiable or differentiable but useless, such
as rounding, median filtering, and compression. The second family, stochastic gradients, appears when
the model randomizes something on every call, so that the gradient of one call does not hold for the
next. The third family, vanishing or exploding gradients, appears when the defense suppresses the
gradient during training, so that the gradient becomes too small or too large to use for a step.

All three share one feature: they make the measurement look better without making the model better,
and because the measurement is usually the only evidence that is reported, this error spreads fast.
This chapter does not say that all such defenses are useless; some raise the cost of attack, and that
cost is real. What is wrong is concluding robustness from a broken measurement.

\section{Seven ways to obfuscate the gradient}
\tingkat{intermediate}

The table below summarizes the seven techniques that appear most often, their mechanism family, and
the symptom that ought to make you suspicious. The numbers in the last column are not a diagnosis,
but a trigger for a question.

\begin{center}
\small
\begin{tabular}{@{}p{0.6cm}p{4.6cm}p{2.1cm}p{6.4cm}@{}}
\hline
No & Technique & Family & Visible symptom \\
\hline
1 & Input or output quantization & shattered & zero gradient; attack accuracy equals clean accuracy \\
2 & Non-differentiable preprocessing: median, blur, compression & shattered & zero gradient, or spikes with no pattern \\
3 & Sample filtering and rejection & shattered & discrete decisions; no gradient \\
4 & Random transforms: resize, pad, shift & stochastic & a different gradient on every call \\
5 & Random classifiers and random ensembles & stochastic & one sample points the wrong way \\
6 & Gradient suppressed during training & vanishing & small gradient despite a large loss \\
7 & Multi-step optimization defenses & vanishing & exploding gradient; attack does not converge \\
\hline
\end{tabular}
\end{center}

\index{quantization}
The first technique, quantization, has already been proved above, and its effect is exact: not
approximately zero, but zero. The second technique works the same way but more dirtily. Median
filtering replaces every value with the median of its neighborhood, so the function stays flat in
most places, but the boundaries between its pieces are no longer evenly spaced. Compression such as
JPEG discards high-frequency components; the derivative exists but is unstable, and a finite
difference gradient can change sign when the step is shifted. The symptom is characteristic: the
gradient is not zero, but large-scale garbage.

The third technique is the crudest form of obfuscation. The model refuses to answer for inputs it
considers strange, for example those whose distance to the data center is too large, then returns the
label ``unknown''. For the attacker this is not a bad gradient but no gradient at all: the function
is an indicator function, whose derivative is zero everywhere except on the decision surface. A
defense of this kind also rejects legitimate new data, and legitimate data is the reason the system
was installed.

The fourth and fifth techniques are the stochastic family. In the fourth, the input is randomly
transformed before it enters the model: the image is enlarged, shifted by one pixel, or some of its
features are dropped. In the fifth, what is randomized is the model: one ensemble is chosen at random
on every call, or the prediction is taken from the average of many models. In both, the gradient of
one call is the gradient of one realization, whereas what the attacker faces is the distribution of
realizations. A bad gradient sample is not due to bias but to variance: the sign of a direction that
keeps changing is not a good direction.

\index{random filtering}
The sixth technique attacks from the inside. Training deliberately makes the
gradient small around the data points, for example by penalizing the gradient norm or
by using saturating activations, so that the loss surface is very flat around the examples
and steep elsewhere. An attacker that steps by $\varepsilon$ in a flat region
goes nowhere, and the gradient really is small, but not because the model is right.
The seventh technique is its relative: the input is first cleaned by a multi-step
optimization procedure before being classified, and the gradient of that pipeline easily
explodes.

Figure~\ref{fig:permukaan} shows the shape of the surface that quantization
produces on two coordinates. Every flat plateau there is one step. Inside
one plateau, the loss does not change at all, so an attacker that takes a small
step stops where it is --- the three short arrows in that figure. An attacker that is allowed
to step far enough can cross to the next plateau, and the dashed arrow
shows that crossing.

\begin{figure}[!htbp]
\centering
\begin{tikzpicture}
\begin{axis}[
  width=0.80\linewidth, height=7.2cm,
  view={52}{24},
  xlabel={$u$}, ylabel={$v$}, zlabel={$L$},
  domain=-0.35:0.35, y domain=-0.35:0.35,
  samples=27, samples y=27,
  zmin=0, zmax=0.52,
  tick label style={font=\scriptsize},
  label style={font=\small},
  colormap={peta}{color=(layerteal!10) color=(layerteal!60)},
  colorbar=false, axis line style={layergrey}, tick style={layergrey},
]
\addplot3[surf, shader=faceted, faceted color=layergrey!40, line width=0.15pt]
  {0.1*round(10*((x-0.15)^2+(y-0.10)^2))};
\draw[-latex, accent, line width=1pt] (axis cs:0.02,0.02,0) -- (axis cs:0.10,0.06,0);
\draw[-latex, accent, line width=1pt] (axis cs:0.02,0.02,0) -- (axis cs:-0.02,0.08,0);
\draw[-latex, accent, line width=1pt] (axis cs:0.02,0.02,0) -- (axis cs:0.04,-0.05,0);
\node[circle, fill=accent, inner sep=1.5pt] at (axis cs:0.02,0.02,0) {};
\draw[-latex, layerteal, line width=1.2pt, dashed]
  (axis cs:0.02,0.02,0) -- (axis cs:0.30,0.30,0.10);
\node[font=\footnotesize, accent, anchor=north west] at (axis cs:0.09,0.05,0)
  {$\nabla L=0$};
\node[font=\footnotesize, layerteal, anchor=south west] at (axis cs:0.28,0.28,0.12)
  {BPDA};
\end{axis}
\end{tikzpicture}
\caption{A loss surface made flat by quantization:
$L(u,v)=\frac{1}{10}\operatorname{round}(10\,s)$ with
$s=(u-0.15)^2+(v-0.10)^2$.
Every plateau of height $0.1$ is one step: inside it $\nabla L=0$,
so the three short arrows from the red point do not change the height at all.
The dashed arrow crosses to the next step, and only an attacker willing to step
as far as the distance between boundaries reaches it.}
\label{fig:permukaan}
\end{figure}

\begin{tangan}
Materials: one long ruler, one coin, one friend, one table.

Prop up one end of the ruler with a book so that its surface is tilted, then place the coin
in the middle. Close your eyes and hold the coin with one finger. Your friend taps the
table slowly, irregularly, and your task is to guess in which direction the ruler rises
--- only from the vibration that reaches your finger.

Repeat ten times and record right or wrong. Then ask your friend to tap harder,
and repeat ten times. Then ask him to stop tapping altogether, and repeat again.

What you should see: harder tapping does not make your guesses right. Tapping that stops
makes your guesses pure guesses, about half right, no better than tossing
a coin. The vibration tells you that something is moving; it does not tell you the direction.
That is the difference between a bad gradient and a gradient that does not exist.
\end{tangan}

Those seven techniques produce seven symptoms. The first five are the list of Athalye and
colleagues; the last two we add from the same typology.

\begin{enumerate}
  \item \textbf{The attack is unbounded.} The report allows a perturbation of any size,
        then concludes robustness from an attack that ignores the bound.
  \item \textbf{The attack uses a single random sample.} Against a defense that randomizes,
        the attacker draws one realization then reports its failure. What is measured is
        one throw of the die, not the probability of winning.
  \item \textbf{The attack is only one step.} A one-step attack is easily stopped by
        an uneven loss surface, and the robustness it implies disappears
        as soon as the step is repeated.
  \item \textbf{Increasing the distortion does not increase success.} If the budget is enlarged
        but the success does not move, what is broken is usually the gradient.
  \item \textbf{The gradient is reported as zero or NaN.} This is the most naked symptom: a gradient
        that is entirely zero on a non-constant model signals a non-differentiable
        piece, not a flat surface.
  \item \textbf{Accuracy under attack equals clean accuracy.} Our addition.
        Equality down to several decimal places at several different $\varepsilon$
        is the mark of an attack that does not move.
  \item \textbf{There is no adaptive attack.} Our addition. The defense is tested only
        with attacks designed before that defense existed, so it has not been tested
        by an attacker who knows that he is being tested.
\end{enumerate}

\kamus{obfuscated gradient}{gradien dikaburkan}{the gradient given to the attacker
does not point in the right direction, usually because there is a non-differentiable or
random piece inside the pipeline, not because the model is strong}

\kamus{quantization}{kuantisasi}{replacing every value with the nearest multiple of
a fixed distance; the result is a step function whose derivative is zero almost everywhere}

\kamus{adaptive attack}{serangan adaptif}{an attack designed with knowledge of
the defense being tested, including knowing that its gradient may be broken}

That list of symptoms is useful for rejecting weak claims but not for accepting strong ones: a defense
can pass all seven and still be fragile, and conversely it can trigger one symptom and still raise the
cost of attack.


\section{Breaking through: BPDA, EOT, and reparameterization}
\tingkat{advanced}

If the gradient that is available is wrong, the attacker has three options. Replace the wrong
gradient with the gradient of a function that approximates it; average the gradient over the
randomness that makes it change; or move the work to another space, where the defense obfuscates
nothing. All three have names, and all three work for different reasons.

\index{bpda}
The first option is called \emph{backward pass differentiable approximation}, abbreviated
BPDA. The idea is as simple as it sounds: on the forward pass use the actual function,
including the rounding or the filtering; on the backward pass skip that piece as if
it were the identity. If the wrapped model is $h(x)=f\bigl(g(x)\bigr)$,
the attacker uses

\begin{equation}
\widehat{\nabla}_x h(x) \;=\; \nabla f\bigl(g(x)\bigr),
\end{equation}

that is, the gradient of $f$ evaluated at the point after wrapping, not the gradient of the
composition. For rounding, this means the attacker stops asking ``what is the derivative of
rounding'' and starts asking ``if the numbers are already rounded, in which direction does $f$
rise fastest''.

Why is this approximation reasonable? Because the error is bounded by how far $g$
shifts the point, not by how strange the shape of $g$ is.

\begin{proposisi}
Let $f$ be continuously differentiable with a Lipschitz gradient, that is,
$\lVert \nabla f(u)-\nabla f(v)\rVert \le L\lVert u-v\rVert$ for all $u,v$, and
let $g$ satisfy $\lVert g(x)-x\rVert\le\delta$. Then
\begin{equation}
\bigl\lVert \nabla f\bigl(g(x)\bigr)-\nabla f(x)\bigr\rVert \;\le\; L\,\delta .
\end{equation}
\end{proposisi}

\begin{proof}
Write $v=g(x)-x$ and use the mean value theorem in integral form,
\begin{equation}
\nabla f(x+v)-\nabla f(x) = \int_0^1 H_f(x+tv)\,v\,dt,
\end{equation}
with $H_f$ the Hessian matrix. Because the spectral norm of $H_f$ is bounded by $L$,
\begin{equation}
\bigl\lVert \nabla f(x+v)-\nabla f(x)\bigr\rVert
\le \int_0^1 \lVert H_f(x+tv)\rVert\,\lVert v\rVert\,dt
\le L\lVert v\rVert \le L\delta .
\end{equation}
\end{proof}

This proposition explains when BPDA can be trusted. If rounding only shifts the
input by $q/2$ and $f$ is smooth enough at that scale, the gradient error is small and
the attack step still moves in the right direction. If $q$ is large so that the shift is
comparable to the scale of the curvature of $f$, the bound $L\delta$ becomes large and nothing is
guaranteed. This proposition does not say that $\widehat{\nabla}$ is the gradient of $h$; that
gradient does not exist. It only says that the direction used is close to the correct direction for
$f$ at the point that has already been wrapped.

The second option is \emph{expectation over transformation}, abbreviated EOT. Against
a stochastic defense, the attacker does not face one function but a distribution of functions.
What makes sense to maximize is not the loss on one realization, but the expected
loss over realizations,

\begin{equation}
\overline{L}(x) \;=\; \mathbb{E}_{T\sim\mathcal{T}}\; L\bigl(f_T(x),y\bigr),
\end{equation}

with $T$ denoting the randomness the defense uses, for example a feature mask or a
transformation parameter.

\index{eot}
\begin{lemma}
If $\mathcal{T}$ is finite and $\nabla_x L\bigl(f_T(x),y\bigr)$ is uniformly bounded
in $T$, then
\begin{equation}
\nabla_x \overline{L}(x) \;=\; \mathbb{E}_{T}\;\nabla_x L\bigl(f_T(x),y\bigr).
\end{equation}
Moreover, the average of $m$ independent samples,
$\frac{1}{m}\sum_{i=1}^{m}\nabla_x L\bigl(f_{T_i}(x),y\bigr)$, is an unbiased estimator
whose variance shrinks proportionally to $1/m$.
\end{lemma}

\begin{proof}
Because $\mathcal{T}$ is finite, the expectation is a finite sum, and the derivative may
enter the sum because the derivative is a linear operator. The sample average has the
same expectation as each of its terms, so it is unbiased, and its variance is the variance
of one term divided by $m$ because the samples are independent.
\end{proof}

This lemma answers an objection that comes up often: if the defense is random, is not any
attack just luck? No, as long as the attacker may call the model
repeatedly. Randomness raises the cost, because the attacker needs $m$ calls per step,
but it does not change the direction being pursued; what changes is only the precision of the
estimate. This is the first example of a pattern that recurs throughout the chapter: a defense that
raises the cost is real, whereas one that only makes measurement harder is illusory.

The third option is reparameterization. Instead of attacking $x$, the attacker attacks
the variable that produces $x$ through a transformation that is not obfuscated. If
the defense wraps the model with a function $g$ that has an inverse, the attacker works in the
space $z$ with $x=g^{-1}(z)$ and optimizes $z$. This is useful when $g$ is a
rotation or another linear transformation that does not destroy information, and it is not
useful when $g$ destroys information: no $z$ can recover numbers
that have already been thrown away.

In practice, the three options are used together, plus two habits. First: start
from several random initial points inside the $\varepsilon$ ball, because a loss surface that
is not convex makes a single initial point easy to trap. Second: do not stop at one
value of $\varepsilon$. Robustness measured at a single budget is robustness that is not
measured.

\begin{resep}
Four steps for testing a defense claim, in this order.
\begin{enumerate}
  \item \textbf{Write the budget first.} Fix $\varepsilon$ or the byte limit,
        the number of queries, and the time. These numbers must not change after the first
        result comes out.
  \item \textbf{Attack by transfer, without queries.} Train a substitute model on the data
        you have, build adversarial examples there, then send them to the victim model.
        This step measures whether the defense merely blocks one model.
  \item \textbf{Spend the remaining queries adaptively.} Only now use BPDA or
        EOT on the victim, and record the queries per step, not just the final result.
  \item \textbf{Report both.} Transfer success, adaptive success, and the cost of each,
        always together with clean accuracy.
\end{enumerate}
A warning about order: step 2 must be finished before step 3. If you measure
transfer after seeing the adaptive results, you already know which examples succeed,
and your transfer numbers leak from the future.
\end{resep}

\kamus{BPDA}{backward pass differentiable approximation}{a technique that replaces the gradient
of a non-differentiable piece with the gradient of the function at the point after that piece,
as if the piece were the identity}

\kamus{EOT}{expectation over transformation}{an attack that maximizes the expected
loss over the randomness of the defense, usually by averaging gradients over several
realizations}

\begin{untukdosen}
This part is usually taught in the wrong order, namely the definition of BPDA first.
A better order is the reverse: show a model whose accuracy does not move at all
under attack, ask the students to guess why, then show that its gradient is
zero. Almost all of them will guess ``the model is strong'', and the silence that follows is a
teaching moment in itself.

Two mistakes need to be anticipated. First, concluding that BPDA ``removes''
the defense; what is true is that BPDA restores access to an approximately correct direction, and
its usefulness depends on how large $\delta$ is relative to the scale of the model curvature.
Second, concluding that stochastic defenses are useless because EOT exists; what is true is that
a stochastic defense raises the per-step cost linearly in $m$, so that the defense changes the
attacker's budget, not the direction of the attack. The exercises at the end of the chapter
force students to compute that cost, not merely name it.
\end{untukdosen}

\section{Worked example: one point on a toy model}
\tingkat{intermediate}

\begin{untukanak}
A shopkeeper has an odd habit: every price is rounded to the nearest rupiah
before it is recorded. If a customer quietly shifts a price by one cent,
the shopkeeper's record does not change at all. That customer can conclude that there is no
point in shifting prices, then stop trying. The conclusion is wrong, but
the reasoning is sensible: what he measured really did not move.

Now imagine another customer. She does not ask whether the record changes because of a one-cent
shift; she asks in which direction the record rises fastest, then shifts the
price far enough that the rounding changes. The shopkeeper's record changes, and
she succeeds without breaking any rule. What separates the two is not cleverness,
but the question that was asked.

A model that uses rounding works like that shopkeeper. A program that only tries
small shifts will stop and report that the model is safe; a program that understands
rounding will step further. So when you read a report that a model
cannot be defeated, ask first: did the attacker stop because it lost, or because it
only tried one way?
\end{untukanak}

Before trusting the table, let us compute one point by hand. The model is made as
simple as possible: eight features with all weights equal to one, so the score is the
sum of the eight features. The data come from two normal distributions centered at $-0.5$ and
$+0.5$, with clean accuracy around $92$ percent. The defense is rounding to
multiples of $q=0.05$ on every feature, plus random filtering with drop probability $r$
that multiplies the rest by $1/(1-r)$. The attacker budget is $\varepsilon=0.25$ in the
$\ell_\infty$ norm.

Take one example that is correctly labeled $+1$,

\begin{equation}
x_0 = (0.30;\ 0.20;\ 0.15;\ 0.10;\ 0.05;\ 0.00;\ -0.10;\ -0.08).
\end{equation}

The sum is $0.62$, so the model answers $+1$ with a thin margin. After
rounding, the vector that enters the classifier is
$(0.30;0.20;0.15;0.10;0.05;0.00;-0.10;-0.10)$ and the score is $0.60$.
Note that one coordinate shifts from $-0.08$ to $-0.10$: that is the work of
rounding.

Now compare two attackers. The naive attacker computes the derivative of the actual
function, that is, the derivative of the composition that contains the rounding. Because rounding
is flat around $x_0$, the derivative is zero in all eight coordinates, so its step is
$x_0 + \varepsilon\cdot 0 = x_0$. The score stays $0.60$ and the label stays $+1$.
That attacker's report will read: attack failed, robust accuracy one hundred
percent of clean accuracy.

The second attacker uses BPDA. It does not ask about the derivative of rounding; it
evaluates the gradient of the classifier at the point that has already been rounded, that is, the
vector $(1,1,\dots,1)$, then steps against the label direction by $\varepsilon$ in each
coordinate,

\begin{equation}
x_1 = x_0 - 0.25\cdot(1,1,\dots,1)
= (0.05;\ -0.05;\ -0.10;\ -0.15;\ -0.20;\ -0.25;\ -0.35;\ -0.33).
\end{equation}

After rounding, that vector becomes
\[
(0.05;\ -0.05;\ -0.10;\ -0.15;\ -0.20;\ -0.25;\ -0.35;\ -0.35)
\]
with score $-1.40$. The label flips to $-1$. One step, the cost of one
gradient computation, and the model that was just reported immune turns out to be wrong with a margin
that is not thin.

This example shows something sharper than ``BPDA works''. The second attacker
does not use information that the first does not have: both know the
architecture, the weights, and $q$. What differs is one decision, namely whether to respect the
non-differentiable piece or skip over it. So the ``robustness'' reported by the first attacker
is not a property of the model, but a property of that attacker's notation.

\begin{contoh}
On the same model, with $q=0.05$, $r=0.5$, and a fixed seed, a finite
difference gradient with step $h=10^{-5}$ is zero in $15\,998$ of $16\,000$ entries;
the two remaining entries are exactly $q/h=5\,000$, because their points happen to sit at a step
boundary. The median finite difference gradient norm is zero, whereas the norm of the
BPDA gradient for one mask that keeps four features is $4$. But the direction of a single
mask is still far from the direction of the weights, with a cosine between $0.35$ and $0.94$ over
$200$ masks. The average gradient over $200$ masks gives a vector of norm $2.738$ and
cosine $0.9982$ with respect to the weight vector.
\end{contoh}

Those numbers illustrate three things. First, a shattered gradient is not merely small; it can
also be large and wrong, and two entries that spike will dominate every step
so that the naive attack fails in a way that looks random. Second, BPDA restores the direction
almost perfectly: cosine $0.9982$ means an angle of about three degrees. Third,
the norm of the BPDA gradient for a single mask, namely $4$, differs from its expected norm, $2.738$;
that difference in scale does not hurt an attack that uses the sign, but it does hurt one that uses
the magnitude of the gradient, and that is why EOT is still needed.

\kertas{Recompute the example above without looking at the text. Sum $x_0$, round to
multiples of $0.05$, then sum again; check that the result is $0.60$. Subtract
$0.25$ from each coordinate, round again, sum, and check that the result is
$-1.40$. After that, answer: what is the smallest value of $q$ that makes this one-step BPDA
attack fail to flip the label, and why does that value depend on the initial margin
$0.60$ and not on $\varepsilon$ alone?}

\section{Computational lab: accuracy at every level of obfuscation}
\tingkat{advanced}

\begin{figure}[!htb]
\centering
\begin{tikzpicture}
\begin{axis}[
  width=0.86\linewidth, height=5.6cm,
  xlabel={budget $\varepsilon$ in the $\ell_\infty$ norm},
  ylabel={attack success (\%)},
  xmin=0.03, xmax=0.37, ymin=-0.6, ymax=30,
  xtick={0.05,0.10,0.15,0.20,0.25,0.30,0.35},
  xticklabels={0.05,0.10,0.15,0.20,0.25,0.30,0.35},
  ytick={0,5,10,15,20,25,30},
  tick label style={font=\footnotesize},
  label style={font=\small},
  legend style={font=\scriptsize, at={(0.02,0.98)}, anchor=north west,
                draw=layergrey!35, fill=white, fill opacity=0.92, text opacity=1},
  grid=major, grid style={layergrey!15},
  axis line style={layergrey},
]
\addplot[color=layerindigo, mark=*, mark size=1.5pt, line width=1pt]
  coordinates {(0.05,2.38) (0.10,5.27) (0.15,8.66) (0.20,12.84)
               (0.25,17.40) (0.30,22.32) (0.35,27.71)};
\addlegendentry{no rounding, BPDA attacker}
\addplot[color=accent, mark=square*, mark size=1.4pt, line width=1pt, dashed]
  coordinates {(0.05,2.34) (0.10,5.33) (0.15,8.72) (0.20,12.92)
               (0.25,17.48) (0.30,22.46) (0.35,27.79)};
\addlegendentry{rounding $q=0.05$, BPDA attacker}
\addplot[color=layergrey, mark=triangle*, mark size=1.6pt, line width=0.9pt,
         dotted]
  coordinates {(0.05,0.00) (0.10,0.01) (0.15,0.01) (0.20,0.01)
               (0.25,0.01) (0.30,0.02) (0.35,0.02)};
\addlegendentry{rounding $q=0.05$, naive attacker}
\end{axis}
\end{tikzpicture}
\caption{Attack success against $\varepsilon$ on the eight-feature toy model,
$2\,000$ examples, average of five seeds. At $\varepsilon=0.25$ the three curves give
$17.40$, $17.48$, and $0.01$ percent; at $\varepsilon=0.35$ they give
$27.71$, $27.79$, and $0.02$ percent. The first two curves coincide within the
seed-to-seed deviation, which is $0.35$ to $1.08$ percentage points.}
\label{fig:kurva-epsilon}
\end{figure}

Now we repeat that experiment on two thousand examples, five seeds, and several levels of
obfuscation, with \texttt{numpy} alone. What matters is not the code itself,
but one decision inside it: the function \texttt{fgsm} computes the gradient by finite
differences on the function that has already been wrapped, whereas \texttt{pgd} uses the gradient
of the classifier at the point after wrapping. The results, swept over seven values of
$\varepsilon$, are in Figure~\ref{fig:kurva-epsilon}.

\begin{lab}
\begin{verbatim}
import numpy as np

D, N, EPS = 8, 2000, 0.25
W = np.ones(D)

def data(rng, n, mu=0.5):
    y = 2 * rng.integers(0, 2, n) - 1
    return rng.normal(mu * y[:, None], 1.0, (n, D)), y

def quant(x, q):
    return x if q == 0 else np.round(x / q) * q

def mask(shape, r, rng):
    if r == 0:
        return np.ones(shape)
    return (rng.random(shape) >= r) / (1 - r)

def logit(x, q, m):
    return (quant(x, q) * m) @ W

def fgsm(X, y, q, r, rng):
    m, h = mask(X.shape, r, rng), 1e-5
    G = np.zeros_like(X)
    for j in range(D):
        e = np.zeros(D)
        e[j] = h
        G[:, j] = (logit(X + e, q, m) - logit(X - e, q, m)) / (2 * h)
    return X - EPS * y[:, None] * np.sign(G)

def pgd(X, y, q, r, rng, m_eot=1, steps=40):
    Xa, a = X.copy(), EPS / 8
    for _ in range(steps):
        G = sum(W * mask(X.shape, r, rng) for _ in range(m_eot))
        Xa = np.clip(Xa - a * y[:, None] * np.sign(G / m_eot),
                     X - EPS, X + EPS)
    return Xa

def run(seed, q, r):
    rng = np.random.default_rng(seed)
    X, y = data(rng, N)
    def pred(Z):
        return np.sign(logit(Z, q, mask(Z.shape, r, rng)))
    ok = pred(X) == y
    def succ(Z):
        return 100 * np.mean(pred(Z)[ok] != y[ok])
    return (100 * ok.mean(), succ(fgsm(X, y, q, r, rng)),
            succ(pgd(X, y, q, r, rng)), succ(pgd(X, y, q, r, rng, 20)))

print("  q     r    clean  naive   bpda1   eot")
for q, r in ((0.00, 0.0), (0.01, 0.0), (0.05, 0.0),
             (0.05, 0.5), (0.20, 0.0)):
    M = np.array([run(s, q, r) for s in (1, 2, 3, 4, 5)]).mean(0)
    print("%5.2f %5.2f  %6.2f  %5.2f  %5.2f  %5.2f"
          % (q, r, M[0], M[1], M[2], M[3]))
\end{verbatim}

The output of that program, as it stands:

\begin{verbatim}
  q     r    clean  naive   bpda1   eot
 0.00  0.00   92.02  17.40  17.40  17.40
 0.01  0.00   92.00   0.02  17.40  17.40
 0.05  0.00   91.97   0.01  17.48  17.48
 0.05  0.50   82.87  12.16  25.48  25.67
 0.20  0.00   91.77   0.00  17.65  17.65
\end{verbatim}
\end{lab}

Read that table line by line. The first line is the baseline without obfuscation:
clean accuracy $92.02$ percent, one-step attack success $17.40$ percent.
That number is small because $\varepsilon=0.25$ really is small compared with the spread of the
data, not because the attack is weak.

The second line is the heart of this chapter. With $q=0.01$, clean accuracy barely changes,
$92.00$ percent, but the naive success falls from $17.40$ to $0.02$ percent,
about $870$ times smaller. The attacker that uses BPDA gets $17.40$ percent,
exactly the baseline. That defense changes nothing that an attacker who computes correctly can
observe; it only changes what an attacker who computes incorrectly sees.

The third line repeats the same thing at $q=0.05$: naive $0.01$ percent, BPDA
$17.48$ percent, clean accuracy $91.97$ percent. Raising $q$ fivefold does not
change the BPDA success, but it makes the naive report look better. That is the
fourth symptom from the list in section two, in numerical form.

The fourth line adds random filtering with $r=0.5$, and two things happen at once.
Clean accuracy falls from $91.97$ to $82.87$ percent, because discarding half the features and
multiplying the rest by two raises the variance of the score. BPDA success rises to $25.48$
percent, above the baseline, not because the defense is worse but because the examples
that are still correct after obfuscation have a smaller margin.

The fifth line shows that coarse quantization does not save the model. With
$q=0.20$, comparable to $\varepsilon$, BPDA success is $17.65$ percent, slightly
above the baseline: a point that sits near a step boundary needs only a small push to
cross.

The variance across seeds is small but not zero: at $q=0.05$ with $r=0$ the standard deviation is
about $0.6$ percentage points, and on the line with random filtering about $1.4$ points. The standard
error of a single proportion at $2\,000$ examples and a rate of $0.17$ is about $0.8$ points, so
both numbers are consistent with ordinary sampling.

There are four things this experiment does not measure. First, the model is linear, so the BPDA
gradient is almost perfect and the cosine $0.9982$ is too good for a real model. Second, the
attacker knows $q$ and $r$. Third, only one value of $\varepsilon$ was tested.
Fourth, there is no search over attack hyperparameters. All of that makes the numbers above
a lower bound on attack success, not an upper bound.


The third shortcoming can be closed, and the result is in Figure~\ref{fig:kurva-epsilon}.
The same program is run at seven values of $\varepsilon$ from $0.05$ to
$0.35$, five seeds, without random filtering. The curve of the model without rounding and the curve
of the model with rounding $q=0.05$ almost coincide: at $\varepsilon=0.25$ the two
give $17.40$ and $17.48$ percent, and at $\varepsilon=0.35$ the two give
$27.71$ and $27.79$ percent. The curve of the naive attacker stays below $0.02$ percent
across that whole range. This defense therefore does not shift the curve and does not
break it; it merely adds a third curve that is flat.


\section{Case study: NIDS, and why evasion remains impractical}
\tingkat{advanced}

Up to here we have been dealing with feature vectors that may be changed at will. A network
intrusion detection system, or NIDS, is a different world.
\index{nids}
A NIDS reads a packet stream, extracts features such as flow duration,
packet count, byte count, TCP flags, and inter-arrival times, then decides
whether that flow is an attack. The decision must be made fast, often in milliseconds, for
thousands of flows per second.

Figure~\ref{fig:rantai} puts that chain on one line. What matters is not
the boxes, but the red mark between flow reconstruction and feature extraction.
That is where the attacker works: he does not touch the model, he touches the packets, and the
packets are what determine the feature values that enter the model. Everything to the right of that
mark sees only the consequence; everything to the left is a cause that he can control, within the
protocol limits discussed below.

\begin{figure}[!htbp]
\centering
\begin{tikzpicture}[font=\footnotesize,
  tahap/.style={draw=layergrey!70, fill=layergrey!6, rounded corners=1.5pt,
    align=center, inner sep=3pt, minimum height=1.1cm, text width=1.95cm},
  panah/.style={-latex, layergrey, line width=0.8pt}]
\node[tahap] (n1) at (0,0) {packets\\on the wire};
\node[tahap] (n2) at (2.50,0) {flow\\reconstruction};
\node[tahap] (n3) at (5.00,0) {feature\\extraction};
\node[tahap] (n4) at (7.50,0) {feature\\scaling};
\node[tahap] (n5) at (10.00,0) {NIDS\\model};
\node[tahap] (n6) at (12.50,0) {decision\\threshold};
\draw[panah] (n1) -- (n2);
\draw[panah] (n2) -- (n3);
\draw[panah] (n3) -- (n4);
\draw[panah] (n4) -- (n5);
\draw[panah] (n5) -- (n6);
\node[circle, fill=accent, inner sep=2pt] (sisip) at (3.75,0) {};
\draw[accent, line width=0.8pt] (sisip) -- ++(0,0.62);
\node[accent, anchor=south, font=\scriptsize] at (3.75,0.66)
  {insertion point: packet timing and size};
\node[layergrey, anchor=north, font=\scriptsize, align=center] at (12.5,-0.62)
  {alert or block};
\draw[-latex, layergrey, line width=0.6pt] (n6.south) -- (12.5,-0.5);
\end{tikzpicture}
\caption{The NIDS processing chain, from the wire to the decision. The red mark is one
insertion point that an attacker can use: features such as packet count, byte count, and
inter-arrival times are all determined to the left of that mark, whereas the model on
the right only reads the consequence. A defense that adds obfuscation to the right of the mark
does not touch the cause, and that is the limit discussed in this section.}
\label{fig:rantai}
\end{figure}

\kamus{NIDS}{network intrusion detection system}{a system that monitors network
traffic and decides whether a packet flow is an attack}

\kamus{problem space}{ruang masalah}{the space of real objects that an attacker may change,
for example packets and timing, as distinct from the feature space in which the model works}

The first difficulty is the distance between the two spaces. An attack is usually formulated in
feature space, whereas the attacker sends packets: every feature change must be translated
into a packet change that is still valid under the protocol and still carries out its malicious
function, and that translation does not always exist. Some features cannot even be changed
by the sender, because the packet and byte counts in the reverse direction are determined by the
receiving system.

The second difficulty is the correlation between features. Duration is tied to packet count and
inter-arrival times, and TCP flags are tied to the transport protocol, so an
attacker cannot raise the duration while lowering both inter-arrival times;
discrete features such as protocol number and packet length add further constraints.

The third difficulty is time. An attack that requires ten thousand queries means real
traffic that passes and is logged, and a trace is the first thing an analyst looks at.
A real cost is a real defense; the absence of a gradient is not.

Numbers from the literature make that distance concrete. On the CSE-CIC-IDS2018 data with fifty
filtered features, three models were trained and tested on twenty percent of the data, then
attacked with FGSM, PGD, and BIM. Before the attack all three were almost perfect: accuracy
$0.997$ for ANN, $0.998$ for SVM, and $0.997$ for CNN. On the day of the attack,
with full knowledge, FGSM dropped ANN to $0.756$ with F1 $0.677$, CNN to
$0.842$ with F1 $0.863$, and SVM to $0.150$ with F1 $0.000$. PGD was slightly
stronger on CNN, $0.584$ with F1 $0.689$, and BIM gave $0.688$ with F1
$0.751$.

Those numbers look like a victory for the attacker, and on a static model they are;
but the same experiment continues with a model retrained every day.
With knowledge of a model that is one day old, SVM returns to $0.998$, ANN to
$0.938$, and CNN to $0.957$ with F1 $0.959$, and on the second day the numbers are between
$0.966$ and $0.998$. That retraining does not use adversarial examples, so
the recovery is not the result of an adversarial defense but a consequence of the fact that the
gradient computed on yesterday's model no longer points in the right direction.

This is the NIDS version of the whole chapter: deliberate gradient obfuscation is fragile, because
an attacker can break it with BPDA, whereas obfuscation that comes from model change lasts longer,
since the attacker must update his knowledge, and updating knowledge means gaining access, and
access leaves a trace. The difference is not in the technique, but in who pays the cost.

\index{transferability}
The authors of that review built an attack tree and a taxonomy of practicality problems, then
concluded that only three paths are viable. The first path is blind perturbation:
adding random bytes to UDP packets, or to TCP packets with the PSH flag, while
recomputing the checksum so that the protocol stays valid. This path needs no knowledge of the
model and no queries, so it is immune to retraining, but it has no
direction. The second path is querying through a side channel, which demands continuous network
access and depends on indicators that change with network load and
hardware. The third path is a substitute model trained on data captured from one's own
botnet, which works through transferability but demands training data that is still
fresh.

All three paths are real, and the authors stress that the distance between research and
practice is wide and that the threat model assumptions in many papers are too generous
to the attacker.

\index{evasion}
The authors' own work, which we will call a working note and not a published
result, moves on the other side of the same problem. On one network flow dataset,
a model reaches F1 $0.967037$ from $46\,558$ out of $48\,145$ test examples.
The question is not ``does the attack succeed'', but ``how many features must be
changed for an example to cross the decision boundary while remaining valid under
the protocol''. With a mixed-integer linear program, the answer depends sharply
on the number of features that may be changed: at seven features, not one of fifty
examples could be moved across; at twenty features, thirty-three of fifty could.
The ratio of distance to the boundary between successful and control examples is about $3.28$
for the linear model and $2.13$ for the neural network.

\index{rule of three}
Two more numbers show the caution involved. Of twenty thousand candidates tested for
protocol consistency, none passed; with the rule of three, the upper bound at $95$ percent
for that rate is about $0.019$ percent. On the model with an enlarged safety margin,
the condition takes the form $\varepsilon+40\delta$ with
$\delta=1.245\times10^{-7}$, so that the effective budget is $0.150004981$ and not
$0.15$.

The limits of that working note must be stated together with its results. It uses one dataset,
fifty test examples, and two model families. It verifies the absence of an attack that
meets the conditions, not the discovery of an attack. The number $0$ out of $20\,000$ is a statement
about the candidate generator that was used, not about all possible packets, and
the result says nothing about dynamic NIDSs.

Now the conclusion that must be stated clearly. The toy model in section five does not
prove anything about real NIDSs; it shows that rounding that erases the
gradient produces a misleading report, and that the mechanism can be broken
at a reasonable cost. It does not measure the cost of an attack in packet space. What can be carried
away from this chapter is a habit: every time a defense reports zero success,
ask about its gradient, then ask who pays the cost.

\begin{tantangan}
A paper reports a defense that lowers attack success from $96$
percent to $3$ percent, with clean accuracy that does not change at all. Their attacker
uses one-step FGSM, $\varepsilon=8/255$, the exact gradient, one initial point,
and one seed. The question is: is that robustness real, and what is the cheapest experiment
that most quickly shows that the $3$ percent figure is an artifact?

First hint: look at what is not being paid for. A defense that genuinely
raises the cost almost always pays something, either accuracy, or latency, or
a number of parameters.

Second hint: if the gradient is computed through a non-differentiable piece, the
$3$ percent figure is not attack success, but the residue of examples that happen to lie
near the decision boundary. Repeat with forty-step PGD, BPDA, and three initial points;
then sweep $\varepsilon$ and report the curve.

If you get stuck, go back to section three.
\end{tantangan}

\begin{untukpeneliti}
Three things in this section are often misquoted. First, the categories shattered, stochastic, and
vanishing are the taxonomy of the obfuscated gradients paper; the list of seven techniques in
section two is our pedagogical decomposition. In that paper there are three categories, nine
defenses tested, seven of which depend on an obfuscated gradient, and
six broken completely plus one broken partially.

Second, BPDA is not a general solution. The error bound $L\delta$ holds only when $f$ is smooth
and $\delta$ is small relative to the scale of the curvature; for rounding with $q$ comparable to
$\varepsilon$ that guarantee is empty. Conversely, EOT is not a cure for an attack that points in
the wrong direction: averaging wrong directions still produces a wrong direction.

Third, the NIDS results we quote come from feature-space experiments with full
knowledge and daily retraining, not from packet-space attacks that were actually run. They
support a narrow claim: retraining restores accuracy against attacks that use
the gradient of an old model. They do not support the claim that NIDSs are immune to an attacker who
may change packets.
\end{untukpeneliti}

\section{A checklist before writing that an attack succeeded}
\tingkat{basic}

This section is a summary that can be used directly. Every item is a question
that must be answerable before a claim about an attack or a defense enters the
abstract.

\begin{enumerate}
  \item \textbf{The threat model is written in full.} What does the attacker know: architecture,
        weights, training data, probabilities, or only labels? How many queries may he
        make, and for how long?
  \item \textbf{The attack is bounded by a reasonable budget.} State $\varepsilon$
        and its norm, or the byte and packet limits if the work is in problem space. An attack
        without a bound measures something that nobody faces.
  \item \textbf{The attack is adaptive.}
        \index{adaptive attack}
        The attacker must know the defense being tested, including the possibility that
        its gradient is broken. A defense tested only by attacks from an earlier
        paper has not been tested.
  \item \textbf{The gradient is inspected before it is trusted.} Report the gradient norm, the fraction
        of entries that are zero, and their correlation with the correct direction. A zero gradient on a
        non-constant model is a finding, not a result.
  \item \textbf{Success is reported together with clean accuracy.} A defense that
        lowers clean accuracy from $92$ to $60$ percent must not be compared
        directly with one that keeps $92$ percent.
  \item \textbf{Several seeds, with mean and deviation.} One seed is one
        throw of the die; a difference between two lines that is smaller than their
        deviation is not a difference.
  \item \textbf{A curve, not a single point.} Sweep the value of $\varepsilon$, the number of steps,
        the number of EOT samples, and the number of queries. A good defense shifts the curve;
        an illusory one shifts a single point.
  \item \textbf{The attacker's compute budget is stated:} how many times the model is called,
        how many hours, how many GPUs? An attack that succeeds after a million queries is not
        the same attack as one that succeeds after a hundred.
  \item \textbf{Code and weights are available.} A claim that cannot be repeated cannot be
        checked. If release is impossible, the procedure and hyperparameters must be detailed
        enough to reconstruct.
  \item \textbf{Negative results are reported as negative results.} A defense that fails
        is information; a report that succeeds because its attack was stopped before
        time is an expensive mistake.
\end{enumerate}

\kamus{rule of three}{aturan tiga}{if no event occurs in $n$ independent trials,
the upper bound at $95$ percent for the probability of that event is about $3/n$}

\begin{ringkasan}
\begin{itemize}
  \item An obfuscated gradient blinds the attacker, it does not make the model strong.
        The seven techniques produce the same symptom: the measurement looks better
        without anything getting better.
  \item For a step function such as rounding, the gradient of the composition is zero at almost every
        point, so a one-step attack produces $x'=x$ and ``accuracy under
        attack'' is precisely the clean accuracy.
  \item BPDA replaces the gradient of a non-differentiable piece with the gradient of the
        classifier at the point after the piece. Its error is bounded by $L\delta$: it makes
        sense when the shift is small relative to the scale of the curvature, and it gives no guarantee
        when it is not.
  \item EOT maximizes the expected loss over the randomness of the defense; the average of $m$
        samples is unbiased with variance proportional to $1/m$. Randomness raises the per-step
        cost, it does not change the direction.
  \item On the eight-feature toy model with $\varepsilon=0.25$, rounding
        $q=0.01$ drops naive success from $17.40$ to $0.02$ percent,
        whereas BPDA stays at $17.40$ percent; $q=0.20$ does not hold BPDA back either.
  \item On NIDS, FGSM drops SVM accuracy from $0.998$ to $0.150$, but
        daily retraining restores it to $0.998$. The durable defenses there
        come from the cost of access and protocol constraints, not from gradient obfuscation.
\end{itemize}
\end{ringkasan}

\section{Exercises}
\tingkat{intermediate}

\begin{latihan}
\begin{enumerate}
  \item A paper writes: ``Our defense lowers attack success from
        $98$ to $0$ percent at $\varepsilon=8/255$.'' State three checks
        before accepting that sentence, and what to expect from each check if
        the defense only obfuscates the gradient.
  \item Derive the BPDA error bound for $f(x)=w\cdot x+b$, show that the error is zero, then
        explain why that result does not carry over to a two-layer model with ReLU activations.
  \item On the toy model of section four, determine the smallest value of $q$ that makes one
        BPDA step fail to flip the label of $x_0$, express it as a function of the initial margin and
        $\varepsilon$, then check it with the program.
  \item You have $1\,000$ queries. The defense randomizes a feature mask with $r=0.5$ and
        refuses to answer if more than half the features are dropped. Compute how many PGD steps
        you can run with $m=20$ and with $m=1$, then explain what you gain and what you pay.
  \item Repeat the experiment of section five with a two-layer model with ReLU activations. Report
        clean accuracy, naive success, and BPDA success for three values of $q$,
        then explain why the cosine is no longer close to one.
  \item A NIDS reads fifty features, seven of which are binary flags and three
        protocol numbers. Compute how many features an attacker can actually change, state
        your assumptions, and compare with a $32\times32$ image that has $3\,072$ values.
  \item Explain why daily retraining restores accuracy against attacks that use
        gradients but not against blind perturbation. Connect this with the EOT lemma and
        the definition of an adaptive attack.
  \item Design an experiment that separates two hypotheses: (a) the defense raises the cost of
        attack, (b) the defense only corrupts the measurement. State the metrics, the number of seeds,
        the shape of the curve, and what results would make you accept (a).
\end{enumerate}
\end{latihan}

\peta{Chapter 30 of Athalye et al., \emph{Obfuscated Gradients Give a False Sense of
Security} (2018): section 3 for the three families of obfuscated gradients, section 4 for
BPDA and adaptive attacks, section 6 for the checklist. In Madry et al., \emph{Towards
Deep Learning Models Resistant to Adversarial Attacks} (2018): section 3 for PGD and
random initial points. In elShehaby and Matrawy, \emph{Evasion Adversarial Attacks Remain
Impractical Against ML-based Network Intrusion Detection Systems}, for the attack tree,
the practicality taxonomy, and the retraining experiment. In Chapter 20 of this book for the rule of
three, and in Chapter 27 for FGSM and PGD.}
